% Chương 1. Phụ trách: TV1. Chuyển từ bản Word của TV1 (26/11/2026), đã rút gọn.
\phutrach{TV1}

\subsection{Bối cảnh}
Thương mại điện tử và dịch vụ giao nhận tạo ra lượng vận đơn rất lớn, phải lưu và cập nhật liên tục. Tại Việt Nam, theo Bộ Công Thương, doanh thu thương mại điện tử năm 2025 ước đạt khoảng 31 tỷ USD, tăng 25,5\% so với năm trước~\cite{vn_ecommerce_2025}. Sản lượng bưu chính năm 2024 đạt trên 3,4 tỷ bưu gửi, trong đó dịch vụ phục vụ thương mại điện tử tăng 30\% so với năm 2023~\cite{vn_buu_chinh_2024}.
\ghichu{TV1 kiểm tra lại hai nguồn này. Bản Word ghi ``3,2 tỷ bưu gửi, 2,4 tỷ kiện thương mại điện tử'' theo EBI 2025, nhưng không tìm được nguồn khớp; các số GMV bốn sàn (Metric, YouNet ECI) và Pitney Bowes được bỏ cho gọn.}

Trong hệ thống quản lý vận đơn, cùng một bảng lớn vừa nhận ghi liên tục vừa phục vụ tra cứu theo nhiều thuộc tính. Mỗi vận đơn được tạo kèm số tiền thu hộ (COD), rồi được cập nhật trạng thái nhiều lần qua các chặng. Đồng thời, hệ thống cần tra một đơn theo mã vận đơn (\code{tracking\_code}), lấy các đơn mới nhất của một cửa hàng theo (\code{sender\_id}, \code{created\_at}) và đếm đơn theo kho đích và trạng thái (\code{dest\_hub\_id}, \code{status}). Mỗi kiểu tra cứu cần một chỉ mục phụ, nên chi phí của chỉ mục phụ về dung lượng, độ trễ đọc và thông lượng ghi trở thành câu hỏi trung tâm khi thiết kế.

\subsection{Vấn đề nghiên cứu}
Khi dữ liệu và tải ghi vượt khả năng của một máy chủ RDBMS, một hướng thường được cân nhắc là chuyển sang cơ sở dữ liệu SQL phân tán (NewSQL), vừa giữ mô hình quan hệ vừa mở rộng theo số node. Hai hệ đại diện lưu chỉ mục phụ theo cách khác nhau về bản chất. MySQL InnoDB lưu bảng và từng chỉ mục phụ thành các cây B+Tree trên một máy~\cite{mysql_innodb}. TiDB mã hóa mỗi dòng và mỗi mục chỉ mục thành một cặp Key-Value, chia theo Region trên các node TiKV, lưu bằng LSM-tree trong RocksDB, nhân bản bằng Multi-Raft và commit bằng giao dịch kiểu Percolator~\cite{pingcap_mapping,percolator}.

Câu hỏi của đề tài là: trên cùng bài toán vận đơn, dung lượng, độ trễ đọc và thông lượng ghi qua chỉ mục phụ thay đổi thế nào giữa hai hệ? Thực nghiệm sơ bộ bằng sysbench (đầu Chương~\ref{sec:dung-luong}) cho thấy giả thuyết ban đầu của nhóm không được xác nhận. Đó là động lực để nhóm xây dựng bộ thực nghiệm có kiểm soát trên bảng \code{shipments}, tách riêng từng nguồn chi phí thay vì kết luận chung hệ nào tốt hơn.

\subsection{Câu hỏi nghiên cứu và giả thuyết}
Bảng~\ref{tab:cau-hoi} liệt kê các câu hỏi nghiên cứu. Giả thuyết được đặt trước khi đo và có thể bị số liệu bác bỏ.
\begin{table}[!htp]
\centering\bangnho
\caption{Câu hỏi nghiên cứu, giả thuyết và lý do có thể sai}\label{tab:cau-hoi}
\begin{tabular}{|C{1.2cm}|L{3.6cm}|L{4.7cm}|L{4.7cm}|}
\hline
\dongdau \textbf{Mã} & \textbf{Câu hỏi} & \textbf{Giả thuyết} & \textbf{Lý do có thể sai} \\ \hline
B.1 & Bảng và chỉ mục chiếm bao nhiêu dung lượng? & Mô hình lý thuyết dự đoán được số đo sau khi hiệu chỉnh & Hệ số nén và phiên bản MVCC chưa dọn làm số đo lệch \\ \hline
RQ1 & Covering index giảm độ trễ đọc bao nhiêu? & R nhỏ: lợi ích covering ở TiDB lớn hơn vì chi phí cố định; R lớn: hai hệ có thể gần nhau & Point\_Get rẻ; TiDB gom handle theo lô; MySQL đọc ngẫu nhiên nếu vượt bộ nhớ đệm \\ \hline
RQ2 & Ở độ chọn lọc nào quét bảng nhanh hơn dùng chỉ mục? & Điểm giao nhau của TiDB thấp hơn MySQL & TiDB lấy dòng theo lô; TiKV ít CPU nên quét bảng cũng đắt \\ \hline
RQ3 & Kiểu khóa chính ảnh hưởng thế nào tới chỉ mục phụ? & Khóa chính dạng chuỗi làm chỉ mục phụ lớn hơn trên cả hai hệ; TiDB nonclustered phải tra hai bước khi tìm theo \code{id} & Nén và mã hóa khóa của TiKV có thể làm chênh lệch nhỏ đi \\ \hline
RQ4 & Thêm chỉ mục làm TPS ghi giảm bao nhiêu? & (a) TiDB giảm ít hơn nhờ LSM-tree; (b) TiDB giảm nhiều hơn vì mỗi chỉ mục thêm cặp KV phải prewrite và commit qua Raft & Hai giả thuyết đối lập, do thực nghiệm phân xử \\ \hline
\end{tabular}
\end{table}

\subsection{Mục tiêu và cách đo}
Đề tài không dừng ở việc so sánh độ trễ hay thông lượng tổng mà tìm nguyên nhân kỹ thuật (Bảng~\ref{tab:muc-tieu}). Với chiều đọc, nhóm ép đường truy cập bằng hint, đo các truy vấn đường chuẩn (\code{SELECT 1}, tra theo khóa chính) và khớp số đo với mô hình $t = a + bR$ để tách độ chọn lọc, toán tử truy cập và chi phí gọi KV. Cấu hình bản sao được kiểm tra trước mỗi lượt đo TiDB. Covering index được đánh giá trên cả ba mặt: đọc (RQ1), ghi (RQ4) và dung lượng (B.1).
\begin{table}[!htp]
\centering\bangfont
\caption{Mục tiêu và cách đo}\label{tab:muc-tieu}
\begin{tabular}{|L{3.4cm}|L{5.6cm}|L{5.6cm}|}
\hline
\dongdau \textbf{Nhóm mục tiêu} & \textbf{Mục tiêu} & \textbf{Cách đo} \\ \hline
Tách nguồn chênh lệch đọc & Tách độ chọn lọc, đường truy cập và chi phí gọi KV & Ép chỉ mục bằng hint, đường chuẩn, mô hình $t = a + bR$ \\ \hline
Cấu hình bản sao & Nêu đúng số TiKV, số bản sao, quorum và khả năng chịu lỗi & Kiểm tra trước khi đo TiDB \\ \hline
Covering trên truy vấn vận đơn & Kế hoạch, độ trễ, chi phí ghi, dung lượng & RQ1, RQ4 và B.1 \\ \hline
\end{tabular}
\end{table}

\subsection{Phạm vi và giới hạn}
Đề tài so sánh TiDB v8.5.5 với MySQL 8.4.11, chạy trong Docker trên WSL2 theo hồ sơ tài nguyên chung XS: mỗi hệ dùng tổng cộng 4 CPU; bộ nhớ đệm dữ liệu là 750 MiB ở TiDB (3 TiKV, mỗi node 250 MiB block cache) và 752 MiB ở MySQL (buffer pool). Dữ liệu có hai mức 1 triệu và 3 triệu dòng; mức 10 triệu dòng chỉ được ngoại suy bằng mô hình dung lượng ở Chương~\ref{sec:dung-luong}, không đo. RQ4 chỉ đo ở 1 triệu dòng với 64 luồng; lý do bỏ mức 16 luồng nêu ở Chương~\ref{sec:ghi}.

RQ3 chỉ đo ảnh hưởng của kiểu khóa chính tới dung lượng chỉ mục phụ và tới việc tra cứu, ở 1 triệu dòng trên một máy. Các thí nghiệm còn lại đều giữ cùng lược đồ \code{PRIMARY KEY (id)} trên hai hệ để cô lập chi phí của riêng chỉ mục phụ; vì vậy các biến thể giảm điểm nóng ghi như \code{AUTO\_RANDOM} của TiDB không được thử.

So với đề cương đã đăng ký, nhóm có hai điều chỉnh. Thứ nhất, đề cương dự kiến 1 đến 10 triệu bản ghi; với hồ sơ XS trên laptop 16 GB, nhóm chỉ đo ở 1 và 3 triệu dòng, còn 10 triệu dòng chỉ được tính bằng mô hình dung lượng. Thứ hai, đề cương dự kiến so sánh độ trễ mạng phân tán ở RQ1; theo nhận xét của giảng viên, nhóm tách riêng chi phí gọi KV và không quy chênh lệch cho mạng, vì mọi node nằm trên một máy.

Mọi container của cụm TiDB (1 PD, 1 TiDB server, 3 TiKV) và MySQL chạy luân phiên trên cùng một máy và một ổ đĩa, nên giữa các node không có độ trễ mạng vật lý và chênh lệch hiệu năng không thể quy cho đường truyền. Khối lượng đo chia trên ba máy (Bảng~\ref{tab:ba-may}); số tuyệt đối chỉ được so giữa hai hệ trên cùng một máy, còn giữa các máy chỉ đối chiếu tỉ lệ. Demo tắt node và thêm node ở Chương~\ref{sec:cai-dat} chỉ minh họa cách cụm vận hành, không dùng để kết luận về chi phí chỉ mục.

\subsection{Cấu trúc báo cáo}
Phần 1 gồm Chương~\ref{sec:kien-truc} (tổng quan và kiến trúc TiDB) và Chương~\ref{sec:cai-dat} (cài đặt, cấu hình, vận hành và demo). Phần 2 gồm Chương~\ref{sec:phuong-phap} (vấn đề, cơ chế và phương pháp), Chương~\ref{sec:dung-luong} (dữ liệu, dung lượng và RQ3), Chương~\ref{sec:doc} (chi phí đọc), Chương~\ref{sec:ghi} (chi phí ghi) và Chương~\ref{sec:tong-hop} (tổng hợp). Lệnh, cấu hình và mã nguồn nằm trong kho mã của nhóm, dẫn ở đầu Chương~\ref{sec:cai-dat}.
